%% kapa: exam=1 round=37 year=2026 period=1 form=A subject=real_estate subject_ko=부동산학원론 seq=10 number=90
%% source: Q-Net 감정평가사 기출문제 (cst00309) artlSeq=5258924 file=raw/1008/5258924/2026년 제37회 감정평가사 1차 1교시.pdf
%% conversion: kordoc 4.24.51 -> tools/gen1.py (auto); review=proofread-sonnet; flags=none
\begin{question}{90}
부동산개발사업에 관한 설명으로 옳은 것을 모두 고른 것은?
\begin{framedtext}
ㄱ. 「도시 및 주거환경정비법」상 재건축사업과 「주택법」상 리모델링사업은 주택 단지 중에서 일부의 동을 대상으로 시행할 수 있다. \par
ㄴ. 「도시 및 주거환경정비법」상 재개발사업과 「도시개발법」상 도시개발사업은 환지방식이 가능하다. \par
ㄷ. 「도시 및 주거환경정비법」상 재개발사업과 재건축사업은 관리처분계획수립을 위한 감정평가를 의뢰할 때 조합총회의 의결을 거쳐 1인 이상의 감정평가법인등을 선정할 수 있다. \par
ㄹ. 「도시 및 주거환경정비법」상 재건축사업과 「빈집 및 소규모주택 정비에 관한 특례법」상 소규모재건축사업은 재건축부담금의 부과대상 행위이다. \par
ㅁ. 「도시 및 주거환경정비법」상 재개발사업과 「빈집 및 소규모주택 정비에 관한 특례법」상 자율주택정비사업은 공용수용이 가능하다. \par
\end{framedtext}
\begin{choices}
  \item ㄱ, ㅁ
  \item ㄴ, ㄹ
  \item ㄱ, ㄷ, ㅁ
  \item ㄴ, ㄷ, ㄹ
  \item ㄱ, ㄴ, ㄷ, ㅁ
\end{choices}
\end{question}
\answer{2}
